\scanpage{113}
tel que \deleted{$D$}

a) $\operatorname{supp}D$ ne contient aucune fibre de $R$ dans $X$

b) le \deleted{diviseur} \inserted{$1$-cycle} induit par $D$ sur la fibre générique de $R$ dans $X$ est un diviseur \deleted{Cartier} [condition tjrs satisfaite si $R\longrightarrow X$ est simple ---]. Sous ces conditions, $D$ est un diviseur \deleted{de Cartier}, [transversal aux fibres de $R$].

\uuline{Démonstration} Soit \deleted{$D_R^{(\unsure{1})}=p_1^{-1}(D)$} $R'$ l'ens des $x\in R$ tels que $p_1$ et $p_2$ soient simples en $x$. C'est un ouvert et on voit qu'il induit sur chaque fibre de $p_1$ (ou $p_2$) un ouvert dense. [car $R'=U_1\cap U_2$, $U_1$ induit sur chaque fibre de $p_1$ un ouvert dense, car $p_1$ est séparable, prouvons le point pour $U_2$ ---]. On est ramené
\margin{facile\ldots}
\margin{112}
